\begin{textsample}{POS Dim 1 – vem\_ed\_28 – Score 60.00 – vem\_ed\_28/t150.txt}  \label{ex:f1_pos_003}
PCI en Español com \textbf{conteúdos} da Logística Fatec Bebedouro No ensino \textbf{superior}, o \textbf{intercâmbio} virtual ganha um \textbf{papel} de destaque por \textbf{promover} aprendizagem cultural, linguística e tecnológica. \textbf{É} \textbf{cada} vez mais \textbf{comum} \textbf{que} práticas educacionais mediadas pelas tecnologias contribuam para desenvolver \textbf{competências} \textbf{específicas} e \textbf{essenciais} do Ensino Superior.
No Centro Paula Souza, uma das modalidades de atividade \textbf{colaborativa} \textbf{que} tem ganhado espaço \textbf{são} os Projetos Colaborativos Internacionais (PCIs / Cesu).
Essa metodologia de \textbf{intercâmbio} virtual aproxima professores e alunos de instituições de ensino \textbf{superior} do mundo para discutirem temas da atualidade a \textbf{partir} de \textbf{conteúdos} \textbf{que} envolvem as \textbf{disciplinas} dos cursos tecnológicos.
Em 2024, a Fatec Bebedouro participou da primeira experiência de PCI em parceria com a Uniminuto da Colômbia.
O projeto teve até o \textbf{momento} duas edições, e a temática de interesse \textbf{comum} entre as duas IES \textbf{foi} logística, vinculando as \textbf{disciplinas} de Logística Verde (Fatec Bebedouro) e Gestão da Cadeia de Suprimentos (Uniminuto).
Os grupos mistos de alunos (brasileiros e colombianos), de \textbf{forma} \textbf{colaborativa}, \textbf{analisaram} como a logística \textbf{sustentável} pode contribuir para o \textbf{desenvolvimento} econômico e social para fins da promoção do bem-estar de comunidades \textbf{locais} e globais, estudaram \textbf{princípios} da logística \textbf{sustentável} e sua \textbf{importância} na cadeia de suprimentos moderna e buscaram \textbf{identificar} \textbf{propostas} \textbf{que} contribuam para diminuição do \textbf{impacto} \textbf{ambiental} causado pelos meios de transporte. \textbf{continuação} Essas atividades desenvolvidas por meio da colaboração internacional \textbf{criam} parcerias \textbf{que} potencializam os estudos dos alunos numa \textbf{perspectiva} global, de \textbf{forma} a \textbf{gerar} \textbf{conhecimento} e inovação, \textbf{que} podem \textbf{ser} \textbf{implementados} em seus \textbf{futuros} \textbf{ambientes} de trabalho.
Além disso, a comunicação em \textbf{idioma} estrangeiro, \textbf{que} inicialmente parece \textbf{ser} uma barreira, aos poucos se \textbf{fortalece}.
Especialmente para os alunos fatecanos brasileiros em \textbf{que} o espanhol \textbf{é} componente \textbf{curricular} em Logística, duas horas-aula semanais ao \textbf{longo} de dois semestres.
Ademais, a internacionalização propicia aos professores e estudantes o rompimento de barreiras e a \textbf{integração} da cultura, ensino, pesquisa e a matriz \textbf{curricular} de seus respectivos cursos, de modo \textbf{que} possam \textbf{criar} possibilidades de \textbf{desenvolvimento} dos países e da qualidade de vida das populações. (Stallivieri, 2002).
Propiciar a oportunidade de experiências internacionais abre caminhos para mais competitividade no \textbf{mercado} global.
Esse \textbf{é} um dos \textbf{papéis} “das instituições \textbf{que} buscam o equilíbrio entre as \textbf{expectativas} regionais e nacionais por um lado e os desafios internacionais por outro” (Stallivieri, 2002, local.10).
Quando o \textbf{intercâmbio} virtual acontece entre as \textbf{disciplinas} das áreas técnico-profissionais do curso, o \textbf{idioma} de comunicação pode \textbf{ser} uma barreira.
Por isso, \textbf{é} \textbf{essencial} \textbf{que} o professor de língua estrangeira \textbf{seja} \textbf{colaborativo} para mediar as \textbf{interações}, encorajando e apoiando os alunos na \textbf{troca} de informações relacionadas ao \textbf{conteúdo} em discussão, bem como nas conversas informais \textbf{que} propiciam \textbf{conhecimento} intercultural.
Salomão (2020) considera oportuno \textbf{que} se componham grupos com alunos com maior domínio para atuar como facilitadores sem minimizar o \textbf{papel} dos demais \textbf{que} também contribuem, equilibrando os diferentes \textbf{papéis} na aprendizagem.
Assim, a língua estrangeira vai \textbf{sendo} o motor \textbf{que} possibilita “\textbf{criar} \textbf{formas} de se trabalhar em conjunto na \textbf{construção} de \textbf{conhecimento} ao realizarem uma atividade \textbf{colaborativa} a distância” (Salomão, 2020, p.166), fornecendo andaimes \textbf{que} auxiliam a \textbf{interação} e potencializam as \textbf{habilidades} trazidas pelos alunos.
Outro aspecto importante \textbf{destacado} por Salomão \textbf{é} \textbf{que} o professor de língua estrangeira auxilia o docente da \textbf{disciplina} técnica, nas \textbf{interações} e na \textbf{construção} do projeto.
Estabelecer tal parceria pode \textbf{ser} positivo para a \textbf{produção} de novos \textbf{materiais}, o \textbf{que} \textbf{torna} a aprendizagem mais independente, inclusiva e \textbf{significativa}.
Por se tratar de uma área especializada cuja \textbf{proposta} \textbf{é} discutir dois \textbf{contextos} diferentes (conceitos e aplicações da Logística no \textbf{contexto} brasileiro e colombiano), ainda \textbf{que} houvesse colaboração por \textbf{parte} dos colombianos na comunicação em português, o \textbf{idioma} \textbf{que} se sobressaiu \textbf{foi} a língua espanhola.
Com isso, esse \textbf{idioma} ganhou força na aprendizagem, no caso dos alunos fatecanos brasileiros.
Esse esforço implica um avanço, inclusive no \textbf{desenvolvimento} da língua para fins \textbf{específicos}, pois os alunos passam a ter contato com textos em língua espanhola e a estudar aspectos \textbf{linguísticos} com base em um arcabouço de \textbf{materiais} atrelado a uma especialidade de atuação no \textbf{mercado} de trabalho (Nadin, 2022).
Dessa \textbf{forma}, a aprendizagem se \textbf{torna} mais atrativa e \textbf{significativa} em termos de aplicação na atuação desse aluno em formação. \textbf{Referências} NADIN, O.
L.
Abordagem Terminológico-Discursiva: pelo resgate da Terminologia e da Terminografia no Ensino e na Aprendizagem de Línguas para Fins Específicos no \textbf{contexto} brasileiro.
Trabalhos em Linguística Aplicada, Campinas, v. 61, p. 97-108, jun. 2022.
STALLIVIERI, L.
O \textbf{processo} de internacionalização nas instituições de ensino \textbf{superior}.
Educação Brasileira - Revista do Conselho de Reitores das Universidades, Brasília, v.24, n.48, p. 35-57, 2002.
Disponível em: https://www.ucs.br/site/midia/arquivos/processo\_internacionalizacao.pdf Acesso em 10 mar. 2025.
SALOMÃO, A.
C.
B.
Intercâmbios virtuais e a internacionalização em \textbf{casa}: \textbf{reflexões} e implicações para a Linguística Aplicada.
São Paulo, Estudos Linguísticos. v.49, n.1, p. 152-174, abr. 2020.
Disponível em: https://revistas.gel.org.br/estudos-linguisticos/article/view/2469/1701.
Acesso em 10 mar. 2025.

% matched lemmas: ambiental, ambiente, analisar, cada, casa, colaborativo, competência, comum, conhecimento, construção, contexto, conteúdo, continuação, criar, curricular, desenvolvimento, destacar, disciplina, específico, essencial, expectativa, forma, fortalecer, futuro, gerar, habilidade, identificar, idioma, impacto, implementar, importância, integração, interação, intercâmbio, linguístico, local, longo, material, mercado, momento, papel, parte, partir, perspectiva, princípio, processo, produção, promover, proposta, que, referência, reflexão, ser, significativo, superior, sustentável, tornar, troca
\end{textsample}
